% !TeX program = pdflatex
\documentclass[11pt,a4paper]{amsart}

\newcommand{\notelanguage}{finnish}
% Shared preamble for course notes and exercise sets.

\usepackage[T1]{fontenc}
\usepackage{lmodern}
\providecommand{\notelanguage}{english}
\usepackage[finnish,english]{babel}
\AtBeginDocument{\expandafter\selectlanguage\expandafter{\notelanguage}}

\newcommand{\notesfinnishlanguage}{finnish}
\ifx\notelanguage\notesfinnishlanguage
\newcommand{\notestheoremname}{Lause}
\newcommand{\noteslemmaname}{Lemma}
\newcommand{\notespropositionname}{Propositio}
\newcommand{\notescorollaryname}{Korollaari}
\newcommand{\notesdefinitionname}{Määritelmä}
\newcommand{\notesexamplename}{Esimerkki}
\newcommand{\notesproblemname}{Tehtävä}
\newcommand{\notesexercisename}{Tehtävä}
\newcommand{\notesremarkname}{Huomautus}
\newcommand{\notessolutionname}{Ratkaisu}
\else
\newcommand{\notestheoremname}{Theorem}
\newcommand{\noteslemmaname}{Lemma}
\newcommand{\notespropositionname}{Proposition}
\newcommand{\notescorollaryname}{Corollary}
\newcommand{\notesdefinitionname}{Definition}
\newcommand{\notesexamplename}{Example}
\newcommand{\notesproblemname}{Problem}
\newcommand{\notesexercisename}{Exercise}
\newcommand{\notesremarkname}{Remark}
\newcommand{\notessolutionname}{Solution}
\fi

\usepackage[a4paper,margin=30mm]{geometry}
\usepackage{microtype}
\usepackage{graphicx}
\usepackage{enumitem}

\newcommand{\exerciseheading}[2]{%
  \par\noindent
  {\large\bfseries #1\par}
  \vspace{0.1em}
  \noindent{\large #2\par}
  \vspace{1.5em}
}

\usepackage{amsmath}
\usepackage{amssymb}
\usepackage{amsthm}
\usepackage{mathtools}
\usepackage{aliascnt}

% Row operations beside the affected row, e.g. \rowop{+2R_1}.
\newcommand{\rowop}[1]{\mathrel{\big|}\,#1}

% Linear systems: the argument gives the number of alignment pairs.
\newenvironment{linearsystem}[1]
{\left\{
  \begin{alignedat}{#1}}
    {
  \end{alignedat}\right.}

  \usepackage[hidelinks,pdfusetitle]{hyperref}

  \numberwithin{equation}{section}
  \setlist{itemsep=0.25em,topsep=0.5em}

  \theoremstyle{plain}
  \newtheorem{theorem}{\notestheoremname}[section]
  \newaliascnt{lemma}{theorem}
  \newtheorem{lemma}[lemma]{\noteslemmaname}
  \aliascntresetthe{lemma}
  \newaliascnt{proposition}{theorem}
  \newtheorem{proposition}[proposition]{\notespropositionname}
  \aliascntresetthe{proposition}
  \newaliascnt{corollary}{theorem}
  \newtheorem{corollary}[corollary]{\notescorollaryname}
  \aliascntresetthe{corollary}

  \theoremstyle{definition}
  \newaliascnt{definition}{theorem}
  \newtheorem{definition}[definition]{\notesdefinitionname}
  \aliascntresetthe{definition}
  \newaliascnt{example}{theorem}
  \newtheorem{example}[example]{\notesexamplename}
  \aliascntresetthe{example}
  \newaliascnt{problem}{theorem}
  \newtheorem{problem}[problem]{\notesproblemname}
  \aliascntresetthe{problem}
  \newtheorem{exercise}{\notesexercisename}

  \theoremstyle{remark}
  \newaliascnt{remark}{theorem}
  \newtheorem{remark}[remark]{\notesremarkname}
  \aliascntresetthe{remark}

  \newenvironment{solution}
  {
  \begin{proof}[\notessolutionname]}
    {
  \end{proof}}

  \usepackage[nameinlink,noabbrev]{cleveref}

  \crefname{theorem}{theorem}{theorems}
  \crefname{lemma}{lemma}{lemmas}
  \crefname{proposition}{proposition}{propositions}
  \crefname{corollary}{corollary}{corollaries}
  \crefname{definition}{definition}{definitions}
  \crefname{example}{example}{examples}
  \crefname{problem}{problem}{problems}
  \crefname{exercise}{exercise}{exercises}
  \crefname{remark}{remark}{remarks}

  \ifx\notelanguage\notesfinnishlanguage
  \crefname{theorem}{lause}{lauseet}
  \crefname{lemma}{lemma}{lemmat}
  \crefname{proposition}{propositio}{propositiot}
  \crefname{corollary}{korollaari}{korollaarit}
  \crefname{definition}{määritelmä}{määritelmät}
  \crefname{example}{esimerkki}{esimerkit}
  \crefname{problem}{tehtävä}{tehtävät}
  \crefname{exercise}{tehtävä}{tehtävät}
  \crefname{remark}{huomautus}{huomautukset}
  \fi

  \newcommand{\NN}{\mathbb{N}}
  \newcommand{\NNone}{\mathbb{N}_1}
  \newcommand{\ZZ}{\mathbb{Z}}
  \newcommand{\QQ}{\mathbb{Q}}
  \newcommand{\RR}{\mathbb{R}}
  \newcommand{\CC}{\mathbb{C}}
  \newcommand{\PP}{\mathbb{P}}

  \DeclareMathOperator{\im}{im}
  \DeclareMathOperator{\rank}{rank}
  \DeclareMathOperator{\syt}{syt}
  \DeclareMathOperator{\pyj}{pyj}
  \DeclareMathOperator{\dom}{dom}
  \DeclarePairedDelimiter{\abs}{\lvert}{\rvert}
  \DeclarePairedDelimiter{\norm}{\lVert}{\rVert}
  \DeclarePairedDelimiterXPP{\maxnorm}[1]{}{\lVert}{\rVert}{_\infty}{#1}
  \newcommand{\set}[1]{\left\lbrace#1\right\rbrace}
  \newcommand{\maxset}[1]{\max\set{#1}}
  \newcommand{\minset}[1]{\min\set{#1}}
  \newcommand{\powerset}[1]{\mathcal{P}\!\left(#1\right)}


\title{Johdatus logiikkaan II\\Harjoituksen 4 ratkaisut}
\author{}
\date{}

\begin{document}

\exerciseheading{Johdatus logiikkaan II}{Harjoituksen 4 ratkaisut}

\begin{exercise}
  Merkitään
  \[
    \begin{aligned}
      2\NN &= \text{parillisten luonnollisten lukujen joukko} \\
      < &= \{(a,b) \in \NN^2 : a < b\} \\
      \mid &= \{(a,b) \in \NN^2 : a \text{ jakaa } b\text{:n}\}
    \end{aligned}
  \]
  Miten voit määritellä parittomien luonnollisten lukujen joukon
  \begin{enumerate}[label=(\alph*)]
    \item \(\{P\}\)-mallissa \(\mathcal M_1 = (\NN,2\NN)\),
    \item \(\{R_0,R_1\}\)-mallissa \(\mathcal M_2 = (\NN,<,\mid)\)?
  \end{enumerate}
\end{exercise}

\begin{exercise}
  Mitkä annetuista termeistä ovat vapaita muuttujalle \(x\) kaavassa
  \[
    (R(x,y) \land \exists zR(x,z)) \to \forall x\exists yR(x,y)?
  \]
  \begin{enumerate}[label=(\alph*)]
    \item \(y\)
    \item \(z\)
    \item \(x\)
    \item \(c\)
  \end{enumerate}
\end{exercise}

\begin{exercise}
  Mikä annetuista muuttujista voidaan vaihtaa sidotun muuttujan \(y\) tilalle
  kaavassa \(\exists yR_0(x,y) \lor \exists zR_1(x,z)\) ilman että kaavan
  merkitys muuttuu:
  \begin{enumerate}[label=(\alph*)]
    \item \(x\)
    \item \(y\)
    \item \(z\)
  \end{enumerate}
\end{exercise}

\begin{exercise}
  Onko termi \(t\) vapaa muuttujalle \(x\) kaavassa \(A\), kun
  \begin{enumerate}[label=(\alph*)]
    \item \(t = y\) ja \(A = \exists yR(x,y)\)
    \item \(t = y\) ja \(A = \exists xR(x,y)\)
    \item \(t = x\) ja \(A = \exists yR(x,y)\)
    \item \(t = z\) ja \(A = \exists zP(z) \land R(x,y)\)
    \item \(t = z\) ja \(A = \exists zP(z) \land R(x,z)\)
  \end{enumerate}
  Jos sijoitus on sallittu, mikä on \(A(t/x)\)?
\end{exercise}

\begin{exercise}
  Todista seuraava erikoistapaus sijoituslemmasta suoraan Tarskin
  totuusmääritelmää käyttäen (huom. ei siis induktiota rakenteen suhteen):
  Olkoon \(A\) kaava \(\exists x(P(x) \land R(x,y))\). Valitaan termiksi
  \(t\) muuttuja \(z\). Tällöin seuraavat ehdot ovat ekvivalentteja kaikilla
  \(\mathcal M\) ja \(s\):
  \begin{enumerate}[label=(\alph*)]
    \item \(\mathcal M \models_s A(t/y)\)
    \item \(\mathcal M \models_{s(a/y)} A\), missä
      \(a = t^{\mathcal M}\langle s\rangle\).
  \end{enumerate}
\end{exercise}

\begin{exercise}
  Todista sijoituslemma: Olkoon \(A\) mielivaltainen kaava ja \(t\) termi,
  joka on vapaa muuttujalle \(y\) kaavassa \(A\). Tällöin seuraavat ehdot
  ovat ekvivalentteja kaikilla \(\mathcal M\) ja \(s\):
  \begin{enumerate}[label=(\alph*)]
    \item \(\mathcal M \models_s A(t/y)\)
    \item \(\mathcal M \models_{s(a/y)} A\), missä
      \(a = t^{\mathcal M}\langle s\rangle\).
  \end{enumerate}
  Koska kurssilla ei vielä ole käsitelty funktiosymboleita, voit olettaa,
  ettei niitä esiinny aakkostossa.
\end{exercise}

\begin{exercise}
  Päättele lause \(\forall xR(x,x)\) lauseesta \(\forall x\forall yR(x,y)\).
\end{exercise}

\begin{exercise}
  Päättele lause \(\neg\forall xP(x)\) lauseesta \(\forall x\neg P(x)\).
\end{exercise}

\end{document}
